% kap01.tex — Newtonsk mekanik og variationsprincipper
% Hamilton-Jacobi-kompendiet

\chapter{Newtonsk mekanik og variationsprincipper}
\label{kap:newton-variation}

\section{Hvorfor forlade Newton?}

Newtons anden lov,
\begin{equation}
  \vb{F} = m\vb{a},
\end{equation}
er det mest direkte udtryk for mekanikkens grundlov, vi har. Den er lokal
i tid, den er intuitiv, og den virker for et enkelt punkt uden
tvangsbetingelser. Men så snart vi vender os mod systemer, hvor
bevægelsen er begrænset af tvang — en kugle der ruller på en skinne, en
pendul der hænger i en stiv stang, en gasmolekyle der er låst inde i en
beholder — bliver Newtons formulering upraktisk. Problemet er ikke, at
loven svigter; problemet er, at den kræver, at vi kender \emph{alle}
kræfter, herunder de tvangskræfter (normalkræfter, snortræk, reaktioner
fra skinner), der holder systemet på dets tilladte bane. Disse
tvangskræfter udfører intet arbejde i den retning, bevægelsen faktisk
sker i, og de er i praksis ofte netop det, vi \emph{ikke} kender og
\emph{ikke} er interesserede i.

Kapitlets mål er at motivere og udlede et alternativt udgangspunkt —
D'Alemberts princip — der elegant elimininerer tvangskræfterne fra
beskrivelsen, uden at vi behøver kende deres størrelse. Dette er første
skridt på vejen til Lagrange- og senere Hamilton-formalismen
(kapitel~\ref{kap:lagrange} og \ref{kap:hamilton}), som igen er
forudsætningen for Hamilton-Jacobi-teorien, der er kompendiets egentlige
mål.

Den fysiske pointe, vi skal holde fast i hele vejen igennem, er denne:
\emph{et system med tvangsbetingelser har færre uafhængige
frihedsgrader end antallet af koordinater, vi nemmest kan nedskrive
det i}. Kunsten er at finde en beskrivelse, der kun bruger de
uafhængige frihedsgrader — og det er præcis, hvad generaliserede
koordinater giver os.

\section{Frihedsgrader og generaliserede koordinater}

Betragt et system af $N$ punktpartikler. Uden tvangsbetingelser har
systemet $3N$ frihedsgrader — de kartesiske koordinater
$\vb{r}_1,\dots,\vb{r}_N$. Er systemet underlagt $k$ uafhængige
tvangsbetingelser, reduceres antallet af uafhængige frihedsgrader til
\begin{equation}
  n = 3N - k.
\end{equation}

I stedet for at arbejde med de $3N$ kartesiske koordinater og $k$
tvangsligninger, indfører vi $n$ \emph{generaliserede
koordinater}\index{generaliserede koordinater}
$q_1,\dots,q_n$, valgt så tvangsbetingelserne er automatisk opfyldt for
alle værdier af $q_i$. De kartesiske koordinater kan da skrives som
funktioner af de generaliserede koordinater (og eventuelt tiden, hvis
tvangsbetingelserne selv er tidsafhængige):
\begin{equation}
  \vb{r}_\alpha = \vb{r}_\alpha(q_1,\dots,q_n,t), \qquad \alpha = 1,\dots,N.
  \label{eq:transformationsligning}
\end{equation}

\begin{eksempel}
For et simpelt plant pendul af længde $\ell$, der svinger i en fast
lodret plan, er den naturlige kartesiske beskrivelse $(x,y)$ med
tvangsbetingelsen $x^2+y^2=\ell^2$: to koordinater, én
tvangsbetingelse, altså $n = 2-1=1$ frihedsgrad. Vi vælger den
generaliserede koordinat $q_1=\theta$, udslagsvinklen fra lodret, og
skriver
\begin{equation}
  x = \ell\sin\theta, \qquad y = -\ell\cos\theta.
\end{equation}
Tvangsbetingelsen $x^2+y^2=\ell^2$ er nu automatisk opfyldt for
\emph{alle} værdier af $\theta$ — vi behøver aldrig nedskrive den igen.
\end{eksempel}

Denne idé — at vælge koordinater, hvori tvangsbetingelserne er ``bygget
ind'' — er selve nøglen til alt, hvad der følger i dette kompendium.

\subsection{Holonome og ikke-holonome tvangsbetingelser}

En tvangsbetingelse, der kan skrives som en ligning i koordinaterne (og
eventuelt tiden),
\begin{equation}
  f(\vb{r}_1,\dots,\vb{r}_N,t) = 0,
\end{equation}
kaldes \emph{holonom}\index{tvangsbetingelse!holonom}\index{holonom tvangsbetingelse|see{tvangsbetingelse, holonom}}. Pendulets $x^2+y^2-\ell^2=0$ er et eksempel.
Holonome tvangsbetingelser er netop dem, der kan elimineres ved et
koordinatvalg som i \eqref{eq:transformationsligning}: de reducerer
antallet af \emph{koordinater}, vi behøver.

En tvangsbetingelse, der ikke kan skrives på denne form — typisk fordi
den kun begrænser hastigheder uden at være en totale-differential af
en funktion af koordinaterne — kaldes \emph{ikke-holonom}\index{tvangsbetingelse!ikke-holonom}. Det klassiske
eksempel er en mønt, der ruller uden at glide på et plan: mønten kan nå
et hvilket som helst punkt på planet med en hvilken som helst
orientering, men rullebetingelsen (ingen glidning) lægger et
øjeblikkeligt bånd på hastighedsvektoren, som ikke integrerer til et
enkelt bånd på positionen.

I dette kompendium beskæftiger vi os \emph{udelukkende} med holonome
tvangsbetingelser. Det er ikke en indskrænkning, der koster os noget
væsentligt: næsten alle systemer, der er relevante for vejen mod
Hamilton-Jacobi-teori og kvantisering, er holonome, og den formelle
maskine (Lagrange- og Hamilton-formalismen), vi bygger op, forudsætter
det.

\section{Virtuel forskydning og virtuelt arbejde}

For at komme fra Newtons love til et princip, der ikke kræver kendskab
til tvangskræfterne, indfører vi begrebet \emph{virtuel
forskydning}\index{virtuel forskydning}.

En virtuel forskydning $\delta\vb{r}_\alpha$ er en tænkt,
infinitesimal ændring af partikel $\alpha$'s position, som

\begin{itemize}
  \item er i overensstemmelse med tvangsbetingelserne på det givne
        tidspunkt $t$ (systemet ``flyttes'' til en anden tilladt
        konfiguration),
  \item sker \emph{momentant} — tiden holdes fast, $\delta t = 0$ — i
        modsætning til en virkelig forskydning $d\vb{r}_\alpha$, som
        sker i løbet af et tidsinterval $dt$, og som derfor kan
        involvere tidsafhængigheden i tvangsbetingelserne selv.
\end{itemize}

Ordet ``virtuel'' markerer netop denne forskel: $\delta\vb{r}_\alpha$ er
ikke en bevægelse, systemet rent faktisk udfører, men en tilladt
\emph{sammenligningskonfiguration}, vi bruger som redskab.

\emph{Virtuelt arbejde}\index{virtuelt arbejde} er det arbejde, en kraft $\vb{F}_\alpha$ ville
udføre under en sådan virtuel forskydning:
\begin{equation}
  \delta W = \sum_{\alpha} \vb{F}_\alpha \cdot \delta\vb{r}_\alpha.
\end{equation}

Den centrale egenskab, vi skal udnytte, er denne: for et system med
\emph{ideale} tvangsbetingelser — dvs.\ tvangsbetingelser, hvis
tvangskræfter $\vb{R}_\alpha$ står vinkelret på enhver tilladt virtuel
forskydning — gælder
\begin{equation}
  \sum_\alpha \vb{R}_\alpha \cdot \delta\vb{r}_\alpha = 0.
  \label{eq:ideal-tvang}
\end{equation}
Dette er ikke en tilfældig antagelse; det er selve definitionen af, at
en tvangsbetingelse er ``glat'' i den forstand, vi har brug for: en
normalkraft fra en glat skinne, spændingen i en stiv, uelastisk snor,
eller reaktionen fra en glat flade opfylder alle \eqref{eq:ideal-tvang},
fordi de netop er rettet vinkelret på de bevægelser, tvangen tillader.
Friktion opfylder den derimod \emph{ikke} i almindelighed — det er
grunden til, at friktion kræver særlig behandling i denne formalisme
(se bemærkningen nedenfor).

\begin{bemaerkning}
Idealiteten \eqref{eq:ideal-tvang} er den fysiske forudsætning, hele
det følgende bygger på. Den er ikke en matematisk trivialitet, men et
udsagn om, hvilke kræfter naturen rent faktisk lader tvangsbetingelser
udøve. At normalkræfter og snortræk opfylder den, skyldes, at de er
\emph{reaktionskræfter}: de justerer sig præcis til den størrelse, der
er nødvendig for at holde tvangen, og ingen mere — de har intet ``eget''
bidrag til bevægelsesenergien langs de tilladte retninger. Friktion er
anderledes, fordi dens størrelse ikke blot bestemmes af tvangen selv,
men af den relative bevægelse; den formalisme, vi bygger her, dækker
derfor konservative systemer og systemer med ideale tvangsbetingelser,
men skal udvides (typisk med Rayleighs dissipationsfunktion) for at
håndtere friktion. Vi kommer ikke til at have brug for dette i resten
af kompendiet.
\end{bemaerkning}

\section{D'Alemberts princip}

Newtons anden lov for partikel $\alpha$ lyder
\begin{equation}
  \vb{F}_\alpha^{(a)} + \vb{R}_\alpha = m_\alpha\ddot{\vb{r}}_\alpha,
\end{equation}
hvor vi har opdelt den samlede kraft i en \emph{anvendt} (påtrykt) kraft
$\vb{F}_\alpha^{(a)}$ — tyngdekraft, fjederkraft, elektrisk kraft, osv.
— og tvangskraften $\vb{R}_\alpha$. Omskrevet:
\begin{equation}
  \vb{F}_\alpha^{(a)} - m_\alpha\ddot{\vb{r}}_\alpha = -\vb{R}_\alpha.
\end{equation}
Vi ganger begge sider skalært med en virtuel forskydning
$\delta\vb{r}_\alpha$ og summerer over alle partikler:
\begin{equation}
  \sum_\alpha \left(\vb{F}_\alpha^{(a)} - m_\alpha\ddot{\vb{r}}_\alpha\right)\cdot\delta\vb{r}_\alpha
  = -\sum_\alpha \vb{R}_\alpha\cdot\delta\vb{r}_\alpha.
\end{equation}
Ifølge \eqref{eq:ideal-tvang} er højresiden nul for ideale
tvangsbetingelser. Vi står tilbage med \emph{D'Alemberts
princip}\index{D'Alemberts princip}:
\begin{equation}
  \boxed{\sum_\alpha \left(\vb{F}_\alpha^{(a)} - m_\alpha\ddot{\vb{r}}_\alpha\right)\cdot\delta\vb{r}_\alpha = 0.}
  \label{eq:dalembert}
\end{equation}

Dette er den afgørende pointe i kapitlet: tvangskræfterne
$\vb{R}_\alpha$ er \emph{helt forsvundet} fra formuleringen. Vi behøver
ikke kende dem for at bruge \eqref{eq:dalembert} — kun de anvendte
kræfter og de tilladte virtuelle forskydninger, som selv er bestemt af
tvangsbetingelsernes geometri.

\begin{bemaerkning}
D'Alemberts princip kaldes ofte ``det generaliserede princip om virtuelt
arbejde'', fordi det udvider den statiske ligevægtsbetingelse
$\sum_\alpha \vb{F}_\alpha^{(a)}\cdot\delta\vb{r}_\alpha = 0$ (velkendt
fra statik: et system er i ligevægt, hvis det virtuelle arbejde af de
anvendte kræfter er nul for alle tilladte virtuelle forskydninger) til
dynamiske systemer, ved at behandle $-m_\alpha\ddot{\vb{r}}_\alpha$ som
en ekstra ``inertikraft''. Denne omtolkning — at accelerationen kan
regnes med som en kraft, der modvirker bevægelsen — er den samme idé,
man møder som ``d'Alembert-kraft'' eller (i mere kavaleret sprogbrug)
``centrifugalkraft'' i roterende referencesystemer.
\end{bemaerkning}

\section{Eksempel: kloden på et glat skråplan}

Lad os se D'Alemberts princip i arbejde på et konkret eksempel, før vi
går videre til den generelle omskrivning i generaliserede koordinater
(kapitel~\ref{kap:lagrange}).

\begin{eksempel}
En kugle med masse $m$ glider uden friktion ned langs et fast skråplan
med hældningsvinkel $\alpha$ (se figur~\ref{fig:skraaplan}). Lad $s$
være afstanden målt langs planet fra toppen. Tvangsbetingelsen er, at
kuglen forbliver på planets overflade; den eneste frihedsgrad er $s$.

Den anvendte kraft er tyngden, $\vb{F}^{(a)} = -mg\,\hat{\vb{y}}$. En
virtuel forskydning langs planet er $\delta\vb{r} = \delta s\,\hat{\vb{s}}$,
hvor $\hat{\vb{s}}$ er enhedsvektoren langs planet, med komponenter
$(\cos\alpha, -\sin\alpha)$ i det viste koordinatsystem.

Kuglens acceleration er $\ddot{\vb{r}} = \ddot{s}\,\hat{\vb{s}}$ (planet
er fast, så bevægelsen er ren translation langs $\hat{\vb{s}}$).
D'Alemberts princip \eqref{eq:dalembert} giver
\begin{equation}
  \left(\vb{F}^{(a)} - m\ddot{\vb{r}}\right)\cdot\delta\vb{r} = 0
  \quad\Longrightarrow\quad
  \left(-mg\,\hat{\vb{y}} - m\ddot{s}\,\hat{\vb{s}}\right)\cdot\hat{\vb{s}}\,\delta s = 0.
\end{equation}
Da $\delta s$ er arbitrær (og ikke-nul), må parentesen selv være nul:
\begin{equation}
  -mg\,\hat{\vb{y}}\cdot\hat{\vb{s}} - m\ddot{s} = 0
  \quad\Longrightarrow\quad
  \ddot{s} = -g\,\hat{\vb{y}}\cdot\hat{\vb{s}} = g\sin\alpha,
\end{equation}
hvor vi har brugt $\hat{\vb{y}}\cdot\hat{\vb{s}} = -\sin\alpha$ ud fra
komponenterne ovenfor. Vi genfinder det velkendte resultat
$\ddot{s}=g\sin\alpha$ — men bemærk, at vi \emph{aldrig} har skullet
opskrive normalkraften $N$ fra planet på kuglen. Den er elimineret
fuldstændigt, fordi den (ideelt) ikke udfører virtuelt arbejde langs
den tilladte bevægelse.
\end{eksempel}

\begin{figure}[htbp]
  \centering
  \begin{tikzpicture}[scale=1.0]
    % Skråplan
    \draw[thick] (0,3) -- (0,0) -- (5,0);
    \draw[thick] (0,3) -- (5,0);
    % Vinkelmarkering
    \draw (1.0,0) arc (0:31:1.0);
    \node at (1.35,0.28) {$\alpha$};
    % Kugle på planet
    \pgfmathsetmacro{\sfrac}{0.45}
    \coordinate (P) at ({5*\sfrac},{3*(1-\sfrac)});
    \filldraw[fill=blue!20,draw=blue!60!black,thick] (P) circle (0.22);
    % s-akse langs planet
    \draw[-{Latex},dashed] (0,3) -- ($(P)+(0.9,-0.54)$);
    \node[rotate=-31] at ($(0,3)+(1.6,-0.97)$) {$s$};
    % Tyngdekraft
    \draw[-{Latex},thick,red] (P) -- ++(0,-1.1) node[below] {$m\vb{g}$};
    % Normalkraft (elimineres i D'Alembert, men vist for reference)
    \draw[-{Latex},thick,green!55!black] (P) -- ++(0.53,0.90) node[above right] {$\vb{R}$ (elimineres)};
  \end{tikzpicture}
  \caption{Kugle på et glat skråplan. Tyngden $m\vb{g}$ er den anvendte
  kraft; tvangskraften (normalkraften) $\vb{R}$ udfører intet virtuelt
  arbejde langs den tilladte bevægelse $\delta s$ og forsvinder derfor
  fra D'Alemberts princip.}
  \label{fig:skraaplan}
\end{figure}

\section{Fra D'Alembert til et variationsprincip}

D'Alemberts princip \eqref{eq:dalembert} er stadig formuleret i
kartesiske koordinater og vektorielt sprog. Det er et vigtigt
mellemtrin, men ikke slutdestinationen. To ting gør det utilfredsstillende
som slutformulering:

\begin{enumerate}
  \item Det er stadig udtrykt via de $3N$ kartesiske koordinater
        $\vb{r}_\alpha$, ikke via de $n$ uafhængige generaliserede
        koordinater $q_i$, som faktisk beskriver systemets frihedsgrader.
  \item Det er et \emph{differentielt} princip: det udtaler sig om
        systemets tilstand ved ét tidspunkt $t$ ad gangen.
\end{enumerate}

Den anden pointe peger frem mod noget dybere. Der findes en anden
klasse af principper i fysikken — \emph{variationsprincipper} — der i
stedet udtaler sig om hele banen $\vb{r}_\alpha(t)$ fra et starttidspunkt
$t_1$ til et sluttidspunkt $t_2$ på én gang, og udpeger den faktiske
bane som den, der gør en bestemt størrelse (et \emph{integral} over
banen) stationær i forhold til alle ``nærliggende'' tænkte baner.

Det klassiske forbillede er Fermats princip i optikken: lysets faktiske
vej mellem to punkter er den, der gør den optiske gangtid stationær
(typisk minimal) i forhold til alle andre tænkte veje. Fermats princip
forudsiger både retlinet forplantning i homogene medier \emph{og}
brydningsloven ved en grænseflade — alt sammen fra ét og samme
udsagn om et integral. Spørgsmålet, vi arbejder os frem mod i
kapitel~\ref{kap:lagrange}, er, om der findes et analogt princip for
mekanisk bevægelse: findes der en størrelse, hvis integral over banen
er stationært netop for den bane, partiklen faktisk følger?

Svaret er ja, og det er netop \emph{Hamiltons
princip}\index{Hamiltons princip}: den faktiske
bane gør \emph{virkningen}\index{virkning (Hamiltons principfunktion)}
\begin{equation}
  S = \int_{t_1}^{t_2} L\,dt
\end{equation}
stationær, hvor $L$ er Lagrange-funktionen\index{Lagrange-funktion}. At D'Alemberts princip —
et differentielt, øjeblik-for-øjeblik-udsagn — faktisk er
\emph{ækvivalent} med et globalt variationsprincip af denne type, er et
af de smukkeste resultater i klassisk mekanik, og det er hovedemnet i
kapitel~\ref{kap:lagrange}, hvor vi udleder Euler-Lagrange-ligningerne
både fra D'Alemberts princip og fra variationsregning direkte.

\begin{bemaerkning}
Det er værd at bemærke allerede her, hvor langt dette fører: netop
fordi Hamiltons princip er formuleret som et udsagn om \emph{hele
banen} snarere end om kræfter ved et enkelt tidspunkt, er det denne
formulering — ikke Newtons love direkte — der generaliserer smukkest
til kvantemekanikken. Feynmans sti-integral-formulering af
kvantemekanikken (som vi vender tilbage til i kapitel~\ref{kap:felter}
og i det senere kompendium om kvantefeltteori) bygger direkte på
virkningsintegralet $S$ og på selve idéen om at summere over
``nærliggende baner''. Variationsprincippet er altså ikke blot en elegant
omskrivning af Newtons mekanik — det er den formulering, der faktisk
overlever overgangen til kvanteteorien.
\end{bemaerkning}

\section{Opsummering}

\begin{itemize}
  \item Newtons love er lokale og kræver kendskab til tvangskræfter,
        som ofte er uinteressante og ukendte.
  \item Holonome tvangsbetingelser kan elimineres ved valg af
        generaliserede koordinater $q_1,\dots,q_n$ ($n$ = systemets
        antal frihedsgrader).
  \item Virtuelle forskydninger $\delta\vb{r}_\alpha$ er momentane,
        tilladte tænkte forskydninger ($\delta t = 0$).
  \item For ideale tvangsbetingelser udfører tvangskræfterne intet
        virtuelt arbejde: $\sum_\alpha \vb{R}_\alpha\cdot\delta\vb{r}_\alpha=0$.
  \item Dette fører til D'Alemberts princip,
        $\sum_\alpha(\vb{F}_\alpha^{(a)}-m_\alpha\ddot{\vb{r}}_\alpha)\cdot\delta\vb{r}_\alpha=0$,
        hvori tvangskræfterne er elimineret.
  \item D'Alemberts princip er stadig kartesisk og differentielt i
        tid. Det næste skridt — kapitel~\ref{kap:lagrange} — er at
        omskrive det i generaliserede koordinater (Lagrange-formalismen)
        og at vise, at det er ækvivalent med et globalt
        variationsprincip (Hamiltons princip), som er den formulering,
        der peger frem mod Hamilton-Jacobi-teori og senere
        kvantisering.
\end{itemize}

\begin{opgave}
Genudled resultatet $\ddot{s}=g\sin\alpha$ fra eksemplet i
afsnit~\ref{eq:dalembert}'s eksempel ved direkte brug af Newtons anden
lov med en eksplicit normalkraft $N$, og overbevis dig om, at $N$ falder
ud af regningen på samme måde, som den gjorde i D'Alembert-formuleringen
— blot ved en ekstra linje udregning i stedet for slet ikke at optræde.
\end{opgave}

\begin{opgave}
Betragt en perle af masse $m$, der er tvunget til at bevæge sig på en
fast, glat, lodret cirkulær ring med radius $R$, under indflydelse af
tyngden. Indfør vinklen $\theta$ målt fra ringens laveste punkt som
generaliseret koordinat, opskriv den tilladte virtuelle forskydning
$\delta\vb{r}$ i retning af $\theta$, og brug D'Alemberts princip til
at udlede bevægelsesligningen for $\theta(t)$.
\end{opgave}
